\documentclass[10pt,a4paper]{amsart}
\usepackage[margin=19mm]{geometry}
\usepackage{amsmath,amssymb,mathtools}
\usepackage[scr=rsfs,cal=euler]{mathalfa}
\usepackage{xfrac}
\usepackage[unicode,hidelinks,bookmarks=false]{hyperref}
\newcommand{\ie}{i.e.\ }
\newcommand{\cf}{cf.\ }
\newcommand{\resp}{respectively\ }
\newcommand{\Subsection}{\subsection}
\providecommand{\nobreakdash}{\nobreak\hbox{-}}
\setlength{\emergencystretch}{2em}
\multlinegap0pt
\usepackage{amssymb}
\usepackage{enumerate}
\usepackage{xcolor}
\usepackage[shortlabels]{enumitem}
\newtheorem{theorem}{Theorem}
\newtheorem{lemma}{Lemma}
\title{Optimal control of the three-dimensional viscous Camassa-Holm equations with sparse controls}
\author{Giang Nguyen Hai Ha}
\date{Source: arXiv:2604.02724v1 (CC BY 4.0)}
\begin{document}
\maketitle

\noindent\emph{Source and licence.} Optimal control of the three-dimensional viscous Camassa-Holm equations with sparse controls. Version arXiv:2604.02724v1 (CC BY 4.0), distributed by arXiv under the Creative Commons Attribution 4.0 International licence, http://creativecommons.org/licenses/by/4.0/, as recorded in the arXiv licence metadata for this exact version and shown on the abstract page https://arxiv.org/abs/2604.02724v1. Full licence notice: \texttt{licenses/arxiv-2604.02724-CC-BY-4.0.txt}.\par\smallskip

\noindent\textit{Editorial scope.} Counted unit: the article's theorem theo6.3 (source0.tex lines 1061--1073). The article's complete original proof of it is delivered with it (source0.tex lines 1099--1222, one \texttt{proof} environment of 124 lines, which the article places away from the statement). No mathematics is written, repaired or re-derived: the statement and the proof are the article's own lines, in the article's own notation, and no retained line is substituted or re-flowed.  Retained uncounted as the source-local context the proof needs: lines 527--529; lines 1078--1084.  Nothing was omitted from any retained span, so the package carries no omission record.  No retained line contains a citation, so the wrapper carries no bibliography and no bibliography entry.  No retained line contains a figure environment, an image-inclusion command or drawing code, so no image is delivered and the package declares no asset dependency.  Editorial formatting only, in addition: an AMS \texttt{amsart} preamble that restates the article's own declarations and macros used by the retained lines, namely the environments \texttt{theorem}, \texttt{lemma} on separate counters, together with the standard packages those lines need.  The retained environments print in this document as Theorem 1, Lemma 1 in order of appearance, whereas the article prints the same items with the numbering of its own full document.  The complete native source of the article as distributed in the arXiv e-print for this exact version is bundled unchanged as \texttt{sources/197755/source0.tex}.
		\begin{equation} \label{4.3}
			\int_{0}^{T} \left( \overline{\lambda}(t)+\gamma \overline{u}(t),u(t)-\overline{u}(t) \right)\,dt \ge 0, \forall u \in U_{a,b},
		\end{equation}

	\begin{lemma}\label{lem6.3}
		There exists $\varepsilon >0$ such that for all $\kappa \in (0,\varepsilon)$, the following inequality holds
		\begin{align}\label{6.25}
			\mathcal{J}''(\overline{u} + \theta (\overline{u}_\kappa - \overline{u} ))( \overline{u}_\kappa - \overline{u})^2 \geq \dfrac{\mu}{2}\| \overline{u}_\kappa - \overline{u} \|^2_{L^2(Q)^3}
		\end{align}
		with $\mu$ is given by $(\textbf{SSC})$. 
	\end{lemma}

	\begin{theorem}\label{theo6.3}
		With the notations above, there exist constants $\varepsilon>0$ and $C>0$ independent of $\overline{u}$ such that it holds
		\begin{enumerate}
			\item Lipchitz stability for $(\textbf{P}_1)$ and $(\textbf{P}_3)$: 
			\begin{equation} \label{lipz_sta_j13}
				\left\Vert \overline{u}_\kappa - \overline{u} \right\Vert_{L^2(Q)^3} \leq C \kappa \ \text{for all} \ 0<\kappa <\varepsilon.
			\end{equation}
			\item H\"older stability for $(\textbf{P}_2)$: 
			\begin{equation} \label{lipz_sta_j2}
				\left\Vert \overline{u}_\kappa - \overline{u} \right\Vert_{L^2(Q)^3} \leq C \kappa^{2/3} \ \text{for all} \ 0<\kappa <\varepsilon.
			\end{equation}
		\end{enumerate}
	\end{theorem}

	\begin{proof}[Proof of Theorem \ref{theo6.3}]
		Since $\overline{u}_\kappa$ is the local solution of $(\textbf{P}_\kappa)$, we have $\mathcal{J}_\kappa(\overline{u}_\kappa)- \mathcal{J}_\kappa(\overline{u}) \leq 0$, which can be written explicitly as
		\begin{equation} \label{6.3_ori}
			\mathcal{J}(\overline{u}_\kappa) - \mathcal{J}(\overline{u}) + \kappa \left(j_i(\overline{u}_{\kappa}) - j_i(\overline{u}) \right) \leq 0, \hspace{20pt} i=1,2,3.
		\end{equation}
		Recalling that \[ \mathcal{J}({\overline{u}_\kappa}) = \int_0^T g(t, y_{\overline{u}_\kappa}(t)) \, dt+ \frac{\gamma}{2} \iint_Q \vert \overline{u}_\kappa(x,t) \vert ^2 \, dxdt, \]
		using Taylor expansion of $\mathcal{J}$ yields
		\begin{equation*}
			\mathcal{J}({\overline{u}_\kappa}) - \mathcal{J}(\overline{u}) = \mathcal{J}'(\overline{u})(\overline{u}_\kappa - \overline{u}) + \frac{1}{2} \mathcal{J}''(\overline{u}+ \theta (\overline{u}_\kappa - \overline{u}))(\overline{u}_\kappa - \overline{u})^2, \theta\in (0,1).
		\end{equation*}
		Since $\overline{u}_\kappa \in U_{a,b}$, the variational inequality \eqref{4.3} and Lemma \ref{lem6.3} imply that
		\begin{equation*}
			\mathcal{J}(\overline{u}_\kappa) - \mathcal{J}(\overline{u}) \geq \frac{\mu}{4} \Vert \overline{u}_k - \overline{u} \Vert_{L^2(Q)^3}^2 .
		\end{equation*}
		Now we consider the difference $j_i(\overline{u})-j_i(\overline{u}_{\kappa})$. \\
		\textbf{Case 1}: $j_i=j_1$.\\
		When $j_i=j_1$, \eqref{6.3_ori} is equivalent to
		\begin{equation*}
			\mathcal{J}(\overline{u}_\kappa) - \mathcal{J}(\overline{u}) + \kappa \iint_Q \left[ \vert \overline{u}_\kappa(x,t) \vert  - \vert \overline{u}(x,t) \vert \right] \, dxdt \leq 0.
		\end{equation*}
		From this we deduce
		\begin{equation*}
			\mathcal{J}(\overline{u}_\kappa) -\mathcal{J}(\overline{u}) \leq \kappa \iint_Q \left[ \vert \overline{u}(x,t) \vert - \vert\overline{u}_\kappa(x,t) \vert \right] \, dxdt \leq \kappa \Vert \overline{u} - \overline{u}_\kappa \Vert_{L^1(Q)^3}.
		\end{equation*}
		Thus
		\begin{equation} \label{6.3.2}
			\frac{\mu}{4} \Vert \overline{u}_k - \overline{u} \Vert_{L^2(Q)^3}^2 \leq \kappa \Vert \overline{u}_\kappa - \overline{u} \Vert_{L^1(Q)^3}.
		\end{equation}
		However, because the domain $Q$ has a finite volume, we get
		\begin{equation} \label{6.3.1}
			\begin{aligned}
				\Vert \overline{u}_\kappa - \overline{u} \Vert_{L^1(Q)^3} &=\iint_Q 1\times \vert \overline{u}_\kappa - \overline{u} \vert \, dxdt \\
				& \leq \left( \iint_Q 1 \, dxdt \right)^{1/2} \left( \iint_Q \vert \overline{u}_\kappa - \overline{u} \vert^2 \, dxdt \right)^{1/2}\\
				&\leq \text{vol}(Q)^{1/2} \Vert \overline{u}_\kappa - \overline{u} \Vert_{L^2(Q)^3}.
			\end{aligned}
		\end{equation}
		Substituting \eqref{6.3.1} into \eqref{6.3.2}, we obtain
		\begin{equation} \label{sta_j1}
			\frac{\mu}{4} \Vert \overline{u}_k - \overline{u} \Vert_{L^2(Q)^3}^2 \leq \kappa \text{vol}(Q)^{1/2} \Vert \overline{u}_\kappa - \overline{u} \Vert_{L^2(Q)^3}.
		\end{equation}
		\textbf{Case 2}: $j_i=j_2$.\\
		Since $\displaystyle j_2(u)=\Vert u\Vert_{L^2(0,T;\mathbb{L}^1(\Omega))}=\left( \int_0^T \Vert u(t) \Vert_{\mathbb{L}^1(\Omega)}^2 \, dt \right)^{1/2}$, \eqref{6.3_ori} can be rewritten as
		\begin{equation} \label{6.3.4}
			\begin{aligned}
				\mathcal{J}(\overline{u}_{\kappa}) - \mathcal{J}(\overline{u}) &\leq \kappa \left[ j_2(\overline{u}) - j_2 (\overline{u}_{\kappa}) \right]\\
				&= \kappa \left( \int_0^T \Vert \overline{u}(t) \Vert_{\mathbb{L}^1(\Omega)}^2 \, dt \right)^{1/2} - \kappa \left( \int_0^T \Vert \overline{u}_{\kappa}(t) \Vert_{\mathbb{L}^1(\Omega)}^2 \, dt \right)^{1/2}.
			\end{aligned}
		\end{equation}
		Notice that $\overline{u}$ is an optimal solution of problem $(\textbf{P})$, meaning that $\mathcal{J}(\overline{u}_{\kappa}) - \mathcal{J}(\overline{u})\geq 0$, we get $j_2(\overline{u}) \geq j_2(\overline{u}_{\kappa})$. Thus we can write
		\begin{align*}
			&\left( \int_0^T \Vert \overline{u}(t) \Vert_{\mathbb{L}^1(\Omega)}^2 \, dt \right)^{1/2} - \left( \int_0^T \Vert \overline{u}_{\kappa}(t) \Vert_{\mathbb{L}^1(\Omega)}^2 \, dt \right)^{1/2} \\
			&\leq \left( \int_0^T \Vert \overline{u}(t) \Vert_{\mathbb{L}^1(\Omega)}^2 \, dt - \int_0^T \Vert \overline{u}_{\kappa}(t) \Vert_{\mathbb{L}^1(\Omega)}^2 \, dt \right)^{1/2} \\
			&= \left[\int_0^T  \left[  \left( \int_{\Omega} |\overline{u}(t)| \, dx \right)^2 - \left(\int_{\Omega} \vert \overline{u}_{\kappa}(t) \vert \, dx \right)^2\right]dt \right]^{1/2}.
		\end{align*}
		Moreover, it follows
		\begin{align*}
			\left( \int_{\Omega} |\overline{u}(t)| \, dx \right)^2 &- \left(\int_{\Omega} \vert \overline{u}_{\kappa}(t) \vert \, dx \right)^2 \\
			&= \left( \int_{\Omega} (\vert \overline{u}(t) \vert + \vert \overline{u}_{\kappa}(t) \vert) \, dx \right) \left( \int_{\Omega} (\vert \overline{u}(t) \vert - \vert \overline{u}_{\kappa}(t) \vert) \, dx \right) \\
			&\leq C(b). \int_{\Omega} \vert \overline{u}(t) - \overline{u}_{\kappa}(t)\vert \, dx
		\end{align*}
		provided that $\Omega$ is a bounded domain.
		Hence, 
		\begin{equation} \label{6.3.5}
			\begin{aligned}
				\int_0^T  \left[  \left( \int_{\Omega} |\overline{u}(t)| \, dx \right)^2 - \left(\int_{\Omega} \vert \overline{u}_{\kappa}(t) \vert \, dx \right)^2\right]dt &\le \int_0^T C(b). \int_{\Omega} \vert \overline{u}(t) - \overline{u}_{\kappa}(t)\vert \, dt \\
				&\le C(b,T).\Vert \overline{u}_{\kappa} - \overline{u}\Vert_{L^1(Q)^3}. 
			\end{aligned}
		\end{equation}
		Substituting \eqref{6.3.5} into \eqref{6.3.4} leads to
		\begin{equation}
			\mathcal{J}(\overline{u}_{\kappa}) - \mathcal{J}(\overline{u}) \leq \kappa C(b,T)^{1/2}.\Vert \overline{u}_{\kappa} - \overline{u}\Vert_{L^1(Q)^3}^{1/2}.
		\end{equation}
		Then by using Lemma \ref{lem6.3} and \eqref{6.3.1}, we obtain
		\begin{align*}
			\frac{\mu}{4} \Vert \overline{u}_k - \overline{u} \Vert_{L^2(Q)^3}^2 &\leq \kappa C(b,T)^{1/2}.\Vert \overline{u}_{\kappa} - \overline{u}\Vert_{L^1(Q)^3}^{1/2} \\
			&\leq \kappa C(b, T)^{1/2} \text{vol}(Q)^{1/4} \Vert \overline{u}_{\kappa_n}- \overline{u} \Vert_{L^2(Q)^3}^{1/2}.
		\end{align*}
		Consequently,
		\begin{equation} \label{sta_j2}
			\Vert \overline{u}_{\kappa} - \overline{u} \Vert_{L^2(Q)^3} \leq \left(\frac{4}{\mu}\right)^{2/3}C(b,T)^{1/3} \text{vol}(Q)^{1/6} \kappa^{2/3}.
		\end{equation}
		\textbf{Case 3}: $j_i=j_3$.\\
		Analogously to Case 2, we also get
		\begin{equation} \label{6.3.6}
			\begin{aligned}
				\mathcal{J}({\overline{u}}_{\kappa}) - \mathcal{J}(\overline{u}) &\leq \kappa \left[ j_3(\overline{u}) - j_3 (\overline{u}_{\kappa}) \right] \\
				&= \kappa\int_\Omega \Vert \overline{u}(x) \Vert_{L^2(0,T)^3}\, dx - \kappa \int_\Omega \Vert \overline{u}_{\kappa} (x)\Vert_{L^2(0,T)^3} \, dx.
			\end{aligned}
		\end{equation}
		The definition of $j_3$ yields
		\begin{equation*}
			\mathcal{J}({\overline{u}}_{\kappa}) - \mathcal{J}(\overline{u}) \leq \kappa \int_\Omega\left[ \left( \int_0^T \vert \overline{u}(x,t) \vert^2 \, dt\right)^{1/2} - \left( \int_0^T \vert \overline{u}_{\kappa}(x,t) \vert^2 \, dt\right)^{1/2} \right]dx.
		\end{equation*}
		On the other hand, we can write
		\begin{align*} 
			&\left( \int_0^T \vert \overline{u}(x,t) \vert^2 \, dt\right)^{1/2} - \left( \int_0^T \vert \overline{u}_{\kappa}(x,t) \vert^2 \, dt\right)^{1/2} \\&= \displaystyle \frac{\int_0^T \vert \overline{u}(x,t) \vert^2 \, dt - \int_0^T \vert \overline{u}_{\kappa}(x,t) \vert^2 \, dt }{\left( \int_0^T \vert \overline{u}(x,t) \vert^2 \, dt\right)^{1/2} + \left( \int_0^T \vert \overline{u}_{\kappa}(x,t) \vert^2 \, dt\right)^{1/2}} \\
			&\leq \frac{\int_0^T (\vert \overline{u}(x,t) \vert^2 - \vert\overline{u}_{\kappa}(x,t)\vert^2) \, dt}{2a\sqrt{T}}.
		\end{align*}
		It follows that
		\begin{align*}
			\int_0^T \left( \vert \overline{u}(x,t) \vert^2 - \vert\overline{u}_{\kappa}(x,t)\vert^2 \right) \, dt &= \int_0^T \left[ (\vert \overline{u}(x,t)\vert+\vert \overline{u}_{\kappa}(x,t)\vert)(\vert \overline{u}(x,t)\vert-\vert \overline{u}_{\kappa}(x,t)\vert) \right] \, dt \\
			&\leq C(T,b)\int_0^T \vert \overline{u}(x,t) - \overline{u}_{\kappa}(x,t) \vert \, dt.
		\end{align*}
		Therefore
		\begin{align*}
			&\left( \int_0^T \vert \overline{u}(x,t) \vert^2 \, dt\right)^{1/2} - \left( \int_0^T \vert \overline{u}_{\kappa}(x,t) \vert^2 \, dt\right)^{1/2} \\
			&\leq \frac{C(T,b)}{2a\sqrt{T}} \int_0^T \vert \overline{u}(x,t) - \overline{u}_{\kappa}(x,t) \vert \, dt.
		\end{align*}
		Finally,
		\begin{align*}
			\mathcal{J}({\overline{u}}_{\kappa}) - \mathcal{J}(\overline{u}) &\leq \kappa \int_\Omega \frac{C(T,b)}{2a\sqrt{T}} \int_0^T \vert \overline{u}(x,t) - \overline{u}_{\kappa}(x,t) \vert \, dt \\
			&\leq \kappa C(T,a,b,\vert\Omega\vert) \Vert \overline{u}_{\kappa} - \overline{u} \Vert_{L^1(Q)^3}.
		\end{align*}
		Applying Lemma \ref{lem6.3} and \eqref{6.3.1} once again, we conclude
		\begin{align*}
			\frac{\mu}{4} \Vert \overline{u}_k - \overline{u} \Vert_{L^2(Q)^3}^2 &\leq \kappa C(T,a,b,\vert\Omega\vert).\Vert \overline{u}_{\kappa} - \overline{u}\Vert_{L^1(Q)^3} \\
			&\leq \kappa C(T,a,b,\vert\Omega\vert) \text{vol}(Q)^{1/2} \Vert \overline{u}_{\kappa_n}- \overline{u} \Vert_{L^2(Q)^3}.
		\end{align*}
		Consequently,
		\begin{equation} \label{sta_j3}
			\Vert \overline{u}_{\kappa} - \overline{u} \Vert_{L^2(Q)^3} \leq \frac{4}{\mu}C(T,a,b,\vert\Omega\vert) \text{vol}(Q)^{1/2} \kappa.
		\end{equation}
		It follows directly from \eqref{sta_j1}, \eqref{sta_j2} and \eqref{sta_j3} that \eqref{lipz_sta_j13} and \eqref{lipz_sta_j2} hold. We achieve the Lipschitz and H\"older convergence of the solutions to the sparse optimal control problems toward the corresponding standard non-sparse ones.
	\end{proof}

\end{document}
